\section*{अध्याय \ref{ch:b3:cytoskeleton-motility} --- कोशिका-कंकाल की गतिकी और कोशिका-गतिशीलता}

\begin{solution}{exo:b3:cytoskeleton-motility:1}
ऐक्टिन: \qty{7}{nm}, गोलिकाकार ऐक्टिन, ध्रुवीय (काँटेदार/नुकीला), ATP;
वल्कुट, उभार, \omterm{def:b3:cytoskeleton-motility:motors}{मायोसिन} के साथ संकुचन। \omterm{def:b3:cytoskeleton-motility:filaments}{सूक्ष्मनलिका}: \qty{25}{nm},
$\alpha\beta$-ट्यूबुलिन द्विलक, ध्रुवीय (धन/ऋण), GTP; दूर तक
परिवहन, तर्कु, \omterm{def:b3:cytoskeleton-motility:cilia}{पक्ष्माभ}। \omterm{def:b3:cytoskeleton-motility:filaments}{मध्यवर्ती तंतु}: \qty{10}{nm},
कुंडलित-कुंडली प्रोटीन (केराटिन, वाइमेंटिन, लैमिन), अध्रुवीय, कोई
न्यूक्लियोटाइड नहीं; कोशिका और केंद्रक की यांत्रिक शक्ति।
\end{solution}

\begin{solution}{exo:b3:cytoskeleton-motility:2}
वह एकलक सांद्रता जिस पर कोई सिरा न बढ़ता है न सिकुड़ता,
$k_{\text{छूट}}/k_{\text{जुड़}}$। दोनों सिरे इसलिए भिन्न हैं कि उपइकाई
जुड़ने के बाद अपना ATP जल-अपघटित कर देती है, इसलिए सिरों पर भिन्न
रासायनिक स्पीशीज़ रहती हैं (काँटेदार सिरे पर आता ATP-ऐक्टिन, नुकीले
सिरे से जाता ADP-ऐक्टिन) जिनकी आसक्तियाँ भिन्न हैं। बराबर \omterm{thm:b3:cytoskeleton-motility:treadmilling}{क्रांतिक
सांद्रताओं} के साथ केवल एक साम्य होता: कोई चक्रचालन नहीं, कोई अभिवाह
नहीं, और दिशात्मक आवागमन की कीमत चुकाने का कोई उपाय नहीं।
\end{solution}

\begin{solution}{exo:b3:cytoskeleton-motility:3}
\omterm{def:b3:cytoskeleton-motility:filaments}{सूक्ष्मनलिकाओं} पर परिधि की ओर: \omterm{def:b3:cytoskeleton-motility:motors}{काइनेसिन}, धन सिरे की ओर। केंद्र की ओर:
कोशिकाद्रव्यी \omterm{def:b3:cytoskeleton-motility:motors}{डाइनीन}, ऋण सिरे की ओर। प्रतिबल-तंतु: \omterm{def:b3:cytoskeleton-motility:motors}{मायोसिन} II,
ऐक्टिन के काँटेदार सिरे की ओर। \omterm{def:b3:cytoskeleton-motility:cilia}{पक्ष्माभ}: \omterm{def:b3:cytoskeleton-motility:cilia}{अक्षसूत्री डाइनीन}, पड़ोसी
युग्मक के ऋण सिरे (यानी आधार) की ओर।
\end{solution}

\begin{solution}{exo:b3:cytoskeleton-motility:4}
उभार (\omterm{def:b3:cytoskeleton-motility:migration}{पर्णपाद} के लिए Rac, तंतुपाद के लिए Cdc42), आसंजन
(\omterm{def:b3:cytoskeleton-motility:migration}{इंटीग्रिन}; Rac और Rho के नीचे), संकुचन (Rho, \omterm{def:b3:cytoskeleton-motility:motors}{मायोसिन} II के ज़रिए),
पिछले भाग का प्रत्याकर्षण (आसंजनों का छूटना, जिसे संकुचन चलाता है)।
\end{solution}

\begin{solution}{exo:b3:cytoskeleton-motility:5}
$C_{c}^{+} = 1/5 = \qty{0.2}{\micro M}$, $C_{c}^{-} = 0.8/1 =
\qty{0.8}{\micro M}$; $C_{\text{स्था}} = (1 + 0.8)/(5 + 1) = \qty{0.3}{\micro
M}$; $J = 5\times 0.3 - 1 = 0.5$ उपइकाई प्रति सेकंड, \qty{1.35}{nm/s}।
\end{solution}

\begin{solution}{exo:b3:cytoskeleton-motility:6}
प्रति सेकंड $800/8 = 100$ ATP। मीटर में $1/(8\times 10^{-7}) =
1.25\times 10^{6}$ s लगते हैं, यानी लगभग दो सप्ताह, और $1.25\times 10^{8}$
डग तथा ATP। विसरण को \num{16000} वर्ष लगते: सौ करोड़ गुना
अंतर।
\end{solution}

\begin{solution}{exo:b3:cytoskeleton-motility:7}
$F_{\max} = 8.3\times 10^{-20}/3.6\times 10^{-8} = \qty{2.3}{pN}$। वही
ऊर्जा किसी लंबे डग पर फैलने से कम बल देती है --- लंबा उत्तोलक।
\omterm{def:b3:cytoskeleton-motility:motors}{काइनेसिन} (छोटा डग, \qty{6}{pN}) भारी भार के लिए ठीक है; \omterm{def:b3:cytoskeleton-motility:motors}{मायोसिन} V (लंबा
डग) प्रति ATP अधिक दूरी तय करती है और तेज़, हलके परिवहन के लिए ठीक है।
\end{solution}

\begin{solution}{exo:b3:cytoskeleton-motility:8}
(क) टैक्सॉल \omterm{def:b3:cytoskeleton-motility:filaments}{सूक्ष्मनलिकाएँ} जमा देता है: तर्कु खोज, पकड़ और सुधार नहीं
कर सकता, समसूत्री विभाजन जाँच-बिंदु पर रुक जाता है; रेंगती कोशिका उभार
तो बनाती रहती है पर दीर्घकालिक ध्रुवता खो देती है। (ख) कोल्चिसीन
\omterm{def:b3:cytoskeleton-motility:filaments}{सूक्ष्मनलिकाएँ} हटा देता है: कोई तर्कु नहीं, समसूत्री विराम; रेंगना
ऐक्टिन पर चलता रहता है पर बिना किसी टिकाऊ दिशा के। (ग) साइटोकैलेसिन
काँटेदार सिरों पर टोपी लगा देता है: कोई उभार नहीं, कोई रेंगना नहीं;
समसूत्री विभाजन चलता है पर कोशिकाद्रव्य-विभाजन विफल हो जाता है, जिससे
द्विकेंद्रकी कोशिकाएँ बनती हैं। (घ) Rac का निरोध: कोई \omterm{def:b3:cytoskeleton-motility:migration}{पर्णपाद} नहीं,
इसलिए कोई रेंगना नहीं; विभाजन अधिकतर अप्रभावित।
\end{solution}

\begin{solution}{exo:b3:cytoskeleton-motility:9}
निवास-काल अर्ध-काल की कोटि का, \qty{20}{s} (काल-नियतांक
$20/\ln 2 \approx \qty{29}{s}$); अचल अंश \qty{10}{\%}।
किनारे पर बहुलकन को आगे बढ़ना और प्रतिगामी बहाव दोनों देने पड़ते हैं:
$1 + 1.5 = \qty{2.5}{\micro m/min}$।
\end{solution}

\begin{solution}{exo:b3:cytoskeleton-motility:10}
कोई बढ़ता सिरा GTP-ट्यूबुलिन की कोई टोपी बनाए रखता है; उसके पीछे
जल-अपघटन तना हुआ GDP-ट्यूबुलिन बनाता है। टोपी का खोना --- कोई यादृच्छिक
घटना जिसकी प्रायिकता वृद्धि धीमी पड़ने पर बढ़ती है --- तनाव मुक्त कर
देता है और सिरा तब तक तेज़ी से सिकुड़ता है जब तक नई टोपी न बन जाए।
\omterm{thm:b3:cytoskeleton-motility:treadmilling}{क्रांतिक सांद्रता} से ठीक ऊपर हर \omterm{def:b3:cytoskeleton-motility:filaments}{सूक्ष्मनलिका} किसी भी क्षण या तो टोपी
पहने बढ़ रही है या बिना टोपी ढह रही है, इसलिए लंबाइयाँ किसी माध्य पर
इकट्ठी होने के बजाय एक लंबी और एक लुप्त होती समष्टि में बँट जाती हैं।
कोशिका को अपने आयतन की तेज़ टोह और चुन-चुनकर स्थिर करने की क्षमता
मिलती है --- जो धन सिरा किसी गुणसूत्रबिंदु या किसी वल्कुटी स्थल तक
पहुँचे वह पकड़ लिया और रख लिया जाता है, बाकी मिनटों में पुनर्चक्रित हो
जाते हैं: खोजो और पकड़ो।
\end{solution}

\begin{solution}{exo:b3:cytoskeleton-motility:11}
इस पैमाने पर जड़त्व नगण्य है और तरल के समीकरण काल-उत्क्रमणीय हैं: जो
गति स्वयं को उलटकर दोहराती है (कोई चप्पू जो बाहर जाकर वापस आए) वह काय
को वहीं लौटा देती है जहाँ से चला था, चाहे हर आघात की चाल कुछ भी हो।
कोई घूमती कुंडली स्वयं को कभी नहीं दोहराती --- उसकी गति काइरल और सतत
है --- इसलिए वह शुद्ध प्रणोद पैदा करती है। \qty{100}{Hz} पर
प्रति चक्कर एक हज़ार प्रोटॉनों के साथ, प्रति मोटर प्रति सेकंड $10^{5}$
प्रोटॉन।
\end{solution}

\begin{solution}{exo:b3:cytoskeleton-motility:12}
वायुमार्ग के \omterm{def:b3:cytoskeleton-motility:cilia}{पक्ष्माभ} हिल नहीं सकते, श्लेष्मा और जीवाणु साफ़ नहीं होते,
और संक्रमण बार-बार होते हैं। उसी \omterm{def:b3:cytoskeleton-motility:cilia}{अक्षसूत्र} पर बने शुक्राणु-कशाभ अचल
रहते हैं: बंध्यता। भ्रूण में गाँठ के चल \omterm{def:b3:cytoskeleton-motility:cilia}{पक्ष्माभ} एक बाईं ओर का बहाव
चलाते हैं जो बाएँ--दाएँ अक्ष तय करता है; बहाव के बिना पक्ष यादृच्छिक रूप
से चुना जाता है, इसलिए आधे रोगियों के अंग उलटे होते हैं।
\end{solution}

\begin{solution}{pb:b3:cytoskeleton-motility:1}
\textbf{1.} $\qty{300}{\micro m^3} = 3\times 10^{-13}$ L; $\times
2\times 10^{-4}$ mol/L $= 6\times 10^{-17}$ mol $= 3.6\times 10^{7}$
अणु; तंतुओं में $1.8\times 10^{7}$।
\textbf{2.} $1.8\times 10^{7}\times \qty{2.7}{nm} = \qty{4.9}{cm}$
तंतु; प्रति तंतु \qty{0.25}{\micro m} की दर से लगभग $2\times 10^{5}$
तंतु।
\textbf{3.} टोपी-प्रोटीन अधिकांश काँटेदार सिरे रोक देता है, और प्रोफ़िलिन
तथा थाइमोसिन-$\beta$4 एकलकों से बँध जाते हैं जिससे सचमुच मुक्त
ATP-ऐक्टिन \omterm{thm:b3:cytoskeleton-motility:treadmilling}{क्रांतिक सांद्रता} के पास रहता है; वृद्धि उन्हीं सिरों तक
सीमित रहती है जिनकी टोपी कोशिका हटाती है।
\textbf{4.} $C_{\text{स्था}} = 2.2/12.9 = \qty{0.17}{\micro M}$; $J = 11.6
\times 0.17 - 1.4 = 0.6$ उपइकाई प्रति सेकंड; \qty{1.6}{nm/s} $=
\qty{0.1}{\micro m/min}$।
\textbf{5.} $1/0.1 = 10$ गुना।
\textbf{6.} \qty{1}{\micro m/min} $= \qty{16.7}{nm/s} = 6.2$ उपइकाई
प्रति सेकंड प्रति तंतु; $\times 2\times 10^{5} = 1.2\times 10^{6}$ ATP
प्रति सेकंड, यानी कोशिका के आवागमन का लगभग $0.1\,\%$।
\textbf{7.} $1/(8\times 10^{-7}) = 1.25\times 10^{6}$ s $\approx 14.5$
दिन; $1.25\times 10^{8}$ डग और ATP।
\textbf{8.} $L^{2}/2D$: एक मीटर के लिए $(1)^{2}/(2\times 10^{-12}) = 5\times 10^{11}$ s
$\approx \num{16000}$ वर्ष; \qty{10}{\micro m} के लिए $(10^{-5})^{2}/(2\times
10^{-12}) = 50$ s। विसरण किसी कोशिका-काय के भीतर काम आता
है; उससे लंबी किसी भी दूरी को मोटर चाहिए।
\textbf{9.} $5\times 10^{4}/6.02\times 10^{23} = 8.3\times 10^{-20}$ J
$= 20\,k_{B}T$।
\textbf{10.} $8.3\times 10^{-20}/8\times 10^{-9} = \qty{10}{pN}$;
दक्षता $6/10 \approx 60\,\%$।
\textbf{11.} $F = 6\pi\times 0.01\times 10^{-7}\times 8\times 10^{-7} =
1.5\times 10^{-14}$ N $= \qty{0.015}{pN}$, जो \omterm{prop:b3:cytoskeleton-motility:physics}{ठहराव बल} से चार सौ गुना
नीचे है: श्यान कर्षण के सामने एक ही मोटर पर्याप्त है; असली भार बाधाएँ
और बंधन हैं।
\textbf{12.} $10^{4}\times 100 = 10^{6}$ ATP प्रति सेकंड, न्यूरॉन के
बजट का $0.1\,\%$: परिवहन सस्ता है। पर ATP गिरते ही मोटर रुक जाती हैं,
इसलिए सूत्रकणिकाओं की विफलता सूत्रयुग्म को उस सब से वंचित कर देती है
जो कोशिका-काय में बनता है --- तंत्रिका-अपह्रासी रोग में देखी जाने वाली
तंत्रिकाक्षी परिवहन की विफलता।
\textbf{13.} \qty{10}{\micro m/min} $= \qty{167}{nm/s}$; $167/2.7 = 62$
उपइकाई प्रति सेकंड प्रति तंतु।
\textbf{14.} $200\times 20 = 4000$ तंतु; $4000\times 62 = 2.5\times
10^{5}$ उपइकाई प्रति सेकंड।
\textbf{15.} $2.5\times 10^{5}$ ATP प्रति सेकंड, बजट का $0.025\,\%$।
\textbf{16.} $\qty{3}{\micro m}/(\qty{10}{\micro m/min}) = \qty{0.3}{min}
= \qty{18}{s}$; जाल में $2.5\times 10^{5}\times 18 = 4.5\times 10^{6}$
उपइकाइयाँ।
\textbf{17.} $1\,\text{pN}\times 2.7\,\text{nm} = \qty{2.7}{pN.nm} =
0.66\,k_{B}T$: उतनी ऊर्जा का कोई उतार-चढ़ाव लगभग प्रायिकता
$e^{-0.66} \approx 0.5$ से होता है --- एक उपइकाई की झिरियाँ लगातार
खुलती हैं, और रैचेट पूरी तरह संभाव्य है।
\textbf{18.} स्वयं बहुलकन धकेलता है, और जीवाणु अपने ही जाल के आगे
``झिल्ली'' का काम करता है; कोशिका के Arp2/3 तंत्र के साथ वह \omterm{def:b3:cytoskeleton-motility:migration}{पर्णपाद} की
चाल और उससे भी आगे पहुँच जाता है --- \qty{0.5}{\micro m/s},
\qty{30}{\micro m/min} तक।
\textbf{19.} $C = K_{d}$ पर $\theta = 1/2$ और $\mathrm{d}\theta/
\mathrm{d}C = 1/(4C)$: $C$ में \qty{2}{\%} का अंतर $\Delta\theta
= 0.005$ देता है; प्रति ओर \num{25000} ग्राहियों के साथ आगे $125$ अधिक
अधिकृत।
\textbf{20.} हर ओर लगभग \num{12500} अधिकृत हैं, जो
$\sqrt{\num{12500}} \approx 110$ से उतार-चढ़ाव करते हैं, इसलिए दोनों ओर
का अंतर लगभग $160$ से उतार-चढ़ाव करता है: किसी भी क्षण $125$ कोलाहल में
खो जाता है। ग्राही लगभग प्रति सेकंड एक बार फिर से बँधते हैं; $N$
सेकंडों पर औसत लेने से कोलाहल $\sqrt{N}$ से घटता है --- दस सेकंड उसे
$50$ तक ले आते हैं और प्रवणता पढ़ ली जाती है।
\textbf{21.} आगे Rac (Cdc42 के साथ), पीछे Rho। हर जगह सक्रिय
Rac के साथ कोशिका चारों ओर उभार बनाती है, फैल जाती है, और कहीं नहीं जाती।
\textbf{22.} $30/3 = \qty{10}{\micro m/min}$: जैसा दिया है।
\textbf{23.} उभार की जगह दाब-चालित बुलबुलन ले लेता है; आसंजन
और संकुचन अब भी ऐक्टिन तथा \omterm{def:b3:cytoskeleton-motility:motors}{मायोसिन} पर निर्भर हैं, इसलिए चक्र के शेष
भाग के लिए ऐक्टिन अनिवार्य बना रहता है।
\textbf{24.} अग्रभाग हर \qty{18}{s} में फिर से बनता है: Rac
क्रियाशीलता का कोई खिसकाव बहुलकन को एक ही जाल-जीवन-काल में नई दिशा दे
देता है, नया अग्रभाग नेतृत्व करता है, और पीछे के धीमे चरण पीछे-पीछे आते हैं।
\textbf{25.} पात्रे चक्रचालन \qty{1.6}{nm/s}; तंत्रिकाक्ष पार करने में
लगभग $14.5$ दिन; उभार के लिए कोई $2.5\times 10^{5}$ ATP प्रति सेकंड।
\end{solution}
